% lec5.tex 
\chapter{Rings}
\lecture{5}{08:06 AM Thu, Oct 23 2025}{} 
$(A, +, \cdot ) $ is a ring if:
\begin{itemize}
  \item [\ding{172} ] $(A, +) $ is an abelian group.
  \item [\ding{173} ] Is associative: $\forall x,y,z \in   A: (xy) z = x(yz) $  .
  \item [\ding{174} ] Is distributive: $\forall x,y,z \in  A: $
    \[
    \begin{cases}
      x(y+z)  &= xy + xz, \\
      (x+y) z &= xz + yz.
    \end{cases}
    \]
  \item [\ding{175} ] The unity element: $\exists e \in   A: x\cdot e = e \cdot x = x.$  
  \item  [\ding{176} ] Is commutative: $\forall x,y \in   A: xy = yx.$ 
\end{itemize}
\begin{example}
 \begin{itemize}
   \item [\ding{172} ]
$\mathbb{Q}, \mathbb{Z}, \RR , \CC $ are rings.
\divider
\item [\ding{174} ] \[
    \mathbb{Z}[i] = \left\{ a+ib: a,b \in   \mathbb{Z}, i^2 =-1 \right\}
\]
\divider
\item [\ding{174} ] $\RR [X]$.
 \end{itemize}
\end{example}
\begin{definition}[]
A subring $B$ of a ring $A$ is a subgroup of $(A, +) $ if 
\[
1 \in  B, \forall x,y \in   B: xy \in  B.
\]
\end{definition}
\begin{definition}[Rings Product]
  Let $A, B$ be two rings. We define the multiplication in $A \times B $ by:
  \[
  \forall (a,b), (c,d)  \in   A \times B: (a,b) (c,d)  = (ac,bd).
  \]
  Then $(A \times B, +, \cdot  ) $  is a ring with unity $(1_{A}, 1_{B}) $. 
\end{definition}
Let $A$ be a ring and $a,b \in  A$.
 \begin{itemize}
   \item[\ding{182} ] We say that $b$ divides $a$ ($a$ is a multiple of $b$) if
     \[
     \exists q \in  A: a = qb .
     \]
   \item [\ding{183} ] $a$ is invertible if $\exists a' \in  A: a a' = 1. $ 
 \end{itemize}
 The set of invertible elements of $A$ is denoted by $A^{*}$. \\
 \noindent
 \textcolor{purple}{
 \underline{
 \ding{46}Remark:
 }
 }\\
 If $B$ is a subgroup of $A$, then an element $A$ of $B$ can be invertible in $A$ but not in $B$ $(2 \in  \mathbb{Z} \subset \mathbb{Q})$. for all $x \in   A:$ 
 \[
 x \cdot  0 = x ( 0 + 0)  = x \cdot  0 + x \cdot  0,
 \]
 thus,
 \[
  0 = x \cdot  0.
 \]
 \begin{proposition}[]
 $(A ^{*}, \cdot ) $ is an abelian group with $1_{A}$ as neutral element.
 \end{proposition}
 \begin{proposition}[]
   \[
     (A \times B ) * = A^{*} \times B^{*} 
   \]
 \end{proposition}
 \begin{proof}
   We have:
   \begin{align*}
     (a,b) \in   (A \times B ) ^{*} & \iff 
     \exists (x,y)  \in   A \times B: 
     (a,b) (x,y)  = (ax, by) = (1_{A}, 1_{B}) .\\
                                    & \iff a \in  A^{*} \wedge  b \in   B^{*}\\
                                    & \iff (a,b) \in  A^{*} \times B^{*} .
   \end{align*}
 \end{proof}
 \begin{definition}[]
   A field $K$ is a ring with $K^{*} = K \backslash \left\{ 0 \right\}$ .
 \end{definition}
 \begin{example}
 \begin{itemize}
   \item [\ding{172} ] $\mathbb{Z}^{*} = \{ \pm 1\} \implies \mathbb{Z}$ is not a field.
   \item[\ding{173} ] $(\RR \left[ X \right]) ^{*} = \RR ^{*}$.
   \item [\ding{174} ] $(\mathbb{Z}[i]) ^{*} = \left\{ \pm 1, \pm i \right\}$.
 \end{itemize}
 \end{example}
 \section{Ring Morphism}
 \begin{definition}[]
 $ f : A \longrightarrow B $ a ring morphism if:
 \[
 \forall x,y \in  A: \begin{cases}
   f(x + y)  &= f(x)  + f(y),  \\
   f(x \cdot  y) &= f(x)  f(y),  \\
   f(1_{B}) &= 1_{B},
 \end{cases}
 \]
 \end{definition}
 \begin{definition}[Ideals]
 An ideal $I \subset A$ is a subgroup $(I, +) $ of $(A, +) $, if
 \[
 \forall x \in  A, \forall a \in  I: xa \in  I.
 \]
 \end{definition}
\noindent \ding{67} \textsc{Notation:} 
$(a)  = a A = \left\{ ab: b \in  A \right\}$ is 
an ideal of $A$ spanned by $a.$ 
\begin{definition}[]
$I$ is a principal ideal 
of $A$ if $\exists  a \in  A$ such that 
$I = (a)  = aA$.
\end{definition}
\begin{example}
If $ f : A \longrightarrow B $ a ring morphism, then 
\[
\text{Ker}  (f)  = 
\left\{ x \in   A: f(x)  = 0_{B} \right\} \text{ is an ideal of $A$ and $Im(f) $ is subring of $B$.}  
\]
\end{example}
\begin{proposition}[]
A ring $A$ is a field if and only if $\left\{ 0 \right\}$ and $A$ are the only ideals of $A$.
\end{proposition}
\begin{proof}
  \[
    ( \implies ) 
  \]
Let $I \neq \left\{ 0 \right\}$ be an ideal of $A$. Then there exists $a \in  I \backslash  \left\{ 0 \right\} 
\subset A^{*}$. So $1 = a ^{-1} a \in  I $. Then for all $x \in   A: x = x \cdot  1 \in  I$, so $I = A.$ 
\[
  ( \impliedby ) 
\]
Let $a \in  A \backslash  \left\{ 0 \right\}$. Then $(a)  = A$, since $0 \neq a \in  (a) $. Then 
$1 = ab, b \in  A$. So $a \in  A^{*}$. So $A^{*} = A \backslash  \left\{ 0 \right\}$ so $A$ is a field.
\end{proof}
\noindent
\textcolor{larratBicep!50!brown}{
\underline{
\ding{46}Remark:
}
}\\
$ f : K \longrightarrow B $ a ring morphism with $K$ is a field. Since $\text{Ker}  (f)  = \left\{ 0 \right\}$ (Ker(f) 
$\neq \left\{ 0 \right\}$), then $f$ is injective, Im(f) is a field.
\section{Quotient Ring}
Let $A$ be a ring and $I$ not included in $A$ 
an ideal, since $(I , +) \rhol (A, +) $, then 
$\frac{A}{I}$ is an abelian group, 
we define $\frac{A}{I}$ the binary operation, by:
\[
\forall \overline{x}, \overline{y}\in  
\frac{A}{I}: \overline{x} \cdot \overline{y} = 
\overline{xy}.
\]
Then binary operation is well defined. Indeed:
\[
\begin{cases}
\overline{x} = \overline{x'} \\
\overline{y} = \overline{y'}
\end{cases}
 \implies 
 \begin{cases}
 x - x' \in  I \\
 y - y' \in  I
 \end{cases}
 \implies 
 \begin{cases}
 x' = x + a, a \in  I \\
 y' = y + b, b \in  I.
 \end{cases}
\]
which implies that
\[
\overline{x'y'} \implies 
\overline{x'} \overline{y'} = 
\overline{x + a} \cdot \overline{y + a }= 
\overline{x}\cdot \overline{y} = \overline{xy}.
\]
\begin{proposition}[]
$(\frac{A}{I}, +, \cdot ) $ is a ring called the quotient ring of $A$ by $I$.
\end{proposition}
\begin{proof}
\[
  (\overline{x} \overline{y})  \overline{z} = 
  (\overline{xy}) \overline{z} = 
  \overline{(xy) z} = 
  \overline{x(yz) } = 
  \overline{x}(\overline{yz}) = \overline{x}
  (\overline{y} \overline{z}) 
\]
\end{proof}
\begin{theorem}[First theorem of isomorphism (ring version)]
Let $ f : A \longrightarrow B $ be a ring morphism. Then
\[
\frac{A}{\text{Ker(f)}  } \simeq \text{Im(f)}  
\]
\end{theorem}
\begin{proof}
  Let 
  \[
  \begin{array}{cccc}
        \overset{\sim}{f}  : &  \frac{A}{\text{Ker}(f)  }  & \longrightarrow & \text{Im}  (f)  \\
  
             &  \overline{x}  & \longmapsto     & \overset{\sim}{f}(\overline{x})  = f(x) . \\ 
  \end{array}
  \]
  \begin{align*}
    \overline{x} =\overline{y} &= 
    x - y \in  \text{Ker(f)}   \\
                               & \iff f(x - y) = 0 
                               \\
                               & \iff  f(x)  = f(y)  \\
                               & \iff \overset{\sim}{f} (\overline{x}) = 
                               \overset{\sim}{f} ()\overline{y}  
                               \iff  \overset{\sim}{f}  \text{ is well defined and injective.}  
  \end{align*}
\end{proof}
% end of file
